\section{Considerações finais}\label{sec:Conclusion}

Este trabalho investigou em que medida modelos estatísticos e de aprendizado de máquina, usando apenas informações disponíveis antes de cada partida, conseguem prever o total de gols e o total de escanteios da Série A do Campeonato Brasileiro. Para isso, as 3.040 partidas de 2017 a 2024 foram reunidas a partir de uma base principal e de cinco fontes de apoio, nem todas independentes da principal, com regras explícitas para divergências e valores ausentes, e as variáveis preditoras foram construídas só com partidas anteriores. A regressão de Poisson com penalização \textit{ridge}, o Random Forest e o XGBoost foram comparados com a média da liga em validação progressiva, com teste de 2019 a 2024, e as escolhas de variáveis e hiperparâmetros foram feitas apenas com as temporadas anteriores a cada temporada de teste.

A resposta à questão de pesquisa é que o histórico pré-jogo permite prever o total de gols um pouco melhor que a média da liga e o total de escanteios praticamente não. Nos gols, o Random Forest reduziu a desviância de Poisson em cerca de 1,7\%, com significância, e o XGBoost ficou próximo dele; a redução resistiu, com significância, a referências mais fortes, a média da temporada anterior e a média corrente da liga. Nos escanteios, a redução de cerca de 0,9\% sobre a média da liga, obtida pela média simples entre o Random Forest e o XGBoost, deixou de ser significativa contra a média corrente da liga, e a significância mudou entre execuções: os modelos acompanham o nível de escanteios de cada temporada, mas não distinguem as partidas entre si com significância. A variação de desempenho entre temporadas superou a diferença entre os modelos, e nenhuma combinação ponderada superou o melhor modelo isolado. Na previsão de passar das linhas, os modelos superaram a frequência histórica para mais de 2,5 gols, mas não para mais de 9,5 escanteios, em que a escolha da binomial negativa, que acomoda a sobredispersão desse alvo, importou mais que a do modelo.

Os dois alvos dependem de informações diferentes. Os gols foram previstos sobretudo pelos índices de força das equipes, os índices GAP e os \textit{pi-ratings}; os escanteios, pela produção de escanteios de cada equipe na temporada e na forma recente. A formação tática da partida-alvo, testada como hipótese para os escanteios, não trouxe ganho. Esses resultados indicam que, com as informações usadas, a maior parte da variação dos totais depende do que acontece durante a partida.

As contribuições do trabalho são uma base integrada e documentada das partidas do Campeonato Brasileiro de 2017 a 2024, com a origem de cada valor e a resolução das divergências entre as fontes, a ser publicada com o código \Rl{https://github.com/PedroVieiraC/previsao-gols-escanteios-brasileirao}; um procedimento de avaliação sem vazamento temporal, com seleção feita apenas no passado de cada temporada e intervalos de confiança por blocos de rodada; e uma medida do quanto se pode prever gols e escanteios antes da partida nesse campeonato, que serve de referência para trabalhos futuros.

As conclusões se limitam a seis temporadas de teste de uma única competição, com diferenças entre modelos da mesma ordem da incerteza das estimativas, sem ablação por grupo de variáveis, e às variáveis disponíveis nas fontes usadas. Como trabalhos futuros, destacam-se: comparar as probabilidades com as cotações do mercado de apostas, o que exige uma fonte com o mercado de mais e menos gols; incorporar desfalques, escalações nominais e gols esperados (\textit{expected goals}), disponíveis no Sofascore a partir de 2022; modelar o placar e os escanteios de cada equipe separadamente, com distribuições para dados sobredispersos; e validar os modelos em outras competições e temporadas.
